% ══════════════════════════════════════════════════════════════════════════════
%  EXPERIMENT 1 — THE ANGLE OF SUNLIGHT                  (Lesson 2, T4 Week 6)
% ══════════════════════════════════════════════════════════════════════════════
\experimentdiv{The Angle of Sunlight}{Measuring a light patch \textbullet{} Repeat trials \textbullet{} A line of best fit \textbullet{} Why we have seasons}{2}\label{sec:1}
\lessoncodes{VC2S8U12 \textbullet{} VC2S8I01, I02, I04, I05, I06}

\vocabhead
\begin{ctab}
\begin{tabular}{|SW{4.2cm}|W{11.6cm}|}
  \hline\hd & \thead{Word} & \thead{Definition (copy from the slide)} \\ \hline
  & \scaf{\textbf{Axial tilt}} & \scaf{Earth's axis leans 23.5\textdegree{} from upright, so each hemisphere leans towards the Sun for half the year.} \\ \hline
  & \textbf{Angle of elevation} & \\ \hline
  & \textbf{Intensity} & \\ \hline
  & \textbf{Solstice} & \\ \hline
  & \textbf{Line of best fit} & \\ \hline
\end{tabular}
\end{ctab}

\begin{lipanel}
\inquiryli{Measure a light patch precisely at several angles, repeat each and draw a line graph with a line of best fit.}
\understandli{Explain why sunlight heats the ground more strongly in summer, when the Sun is higher in the sky.}
\sctitle
\begin{itemize}[leftmargin=1.4em, itemsep=0pt, topsep=1pt]
  \item I can count the squares in a light patch the same way every time.
  \item I can draw a line of best fit and describe the trend.
  \item I can use my results to explain why Melbourne is warmer in summer.
\end{itemize}
\end{lipanel}

\begin{skillsfocus}
\focus{QP1}{Write a testable hypothesis.}
\focus{PC2}{Change only the angle.}
\focus{PC4}{Count the patch precisely, three times.}
\focus{PA1}{Record results in a THUD table.}
\focus{PA2}{Draw a line of best fit.}
\focus{PA5}{Describe the trend with numbers.}
\focus{EV2}{Link the results to the seasons.}
\end{skillsfocus}

% ─── TECHNIQUE ────────────────────────────────────────────────────────────────
\exppart{Technique spotlight}{Blue}

\begin{technique}{Measuring the light patch}{PC4}
\begin{minipage}[c]{0.6\linewidth}
Trace the edge of the bright patch on the grid paper in pencil. Count every square that is \textbf{more than half} inside the line; ignore the rest. Each square is 1\,cm$^2$, so the count is the area in cm$^2$. Use the same rule for every patch, and have a second person count it too.
\end{minipage}\hfill
\begin{minipage}[c]{0.36\linewidth}\centering
\begin{tikzpicture}[scale=0.42]
  \draw[black!30, step=1] (0,0) grid (7,5);
  \draw[Ink, line width=1.4pt] (3.5,2.5) ellipse (2.6 and 1.6);
  \foreach \x/\y in {1.5/1.5,2.5/1.5,3.5/1.5,4.5/1.5,5.5/1.5,1.5/2.5,2.5/2.5,3.5/2.5,4.5/2.5,5.5/2.5,1.5/3.5,2.5/3.5,3.5/3.5,4.5/3.5,5.5/3.5}{\node[font=\tiny] at (\x,\y) {$\bullet$};}
\end{tikzpicture}
\end{minipage}
\end{technique}

\begin{workedex}{Counting and averaging}
A patch was counted three times at 45\textdegree: 34, 36 and 35 squares.\qquad Average $= (34 + 36 + 35) \div 3 = 35$\,cm$^2$.
\end{workedex}

\begin{wedo}{Counting a patch}
Count the marked squares in the patch in the technique box. What is its area?
\answerbox{2cm}
\end{wedo}

\begin{technique}{Line of best fit}{PA2}
Measurements always have small errors, so the points will not sit perfectly on a line. Draw \textbf{one smooth curve} through the middle of the points, with about as many above as below. Do not join the dots. The curve shows the trend and lets you read values between your measurements.
\end{technique}

% ─── PLAN AND SAFETY / METHOD ────────────────────────────────────────────────
\exppart{Plan, safety and method}{Teal}

\textbf{Aim:} To investigate the effect of changing the angle of the light on the area of the light patch.

\task Write a hypothesis. \skr{QP1}
\answerbox{2.5cm}

\needspace{12\baselineskip}
\task Complete the variables table. \skr{PC2}
\begin{ctab}
\begin{tabular}{|SW{4.4cm}|W{10.2cm}|}
  \hline\hd & \thead{Variable} & \thead{In this experiment} \\ \hline
  & \scaf{Independent} & \scaf{Angle between the torch beam and the paper (\textdegree)} \\ \hline
  & Dependent & \\ \hline
  & Controlled & \\ \hline
  & Controlled & \\ \hline
\end{tabular}
\end{ctab}

\begin{center}
\begin{tikzpicture}[scale=0.8]
  \draw[Ink, line width=1.2pt] (0,0) -- (7,0);
  \fill[black!20] (2.3,-0.1) rectangle (4.9,0.0);
  \begin{scope}[shift={(1.2,3.0)}, rotate=-51.3]
    \draw[eq, fill=white] (-1.3,-0.22) rectangle (0,0.22);
    \draw[eq, fill=black!15] (0,-0.34) rectangle (0.25,0.34);
    \coordinate (h1) at (0.25,0.3); \coordinate (h2) at (0.25,-0.3); \coordinate (tb) at (-1.3,0);
  \end{scope}
  \draw[Ink, dashed] (h1) -- (4.9,0) (h2) -- (2.3,0);
  \draw[Ink, densely dotted] (1.2,3.0) -- (3.6,0);
  \draw[Ink, line width=1pt] (2.8,0) arc[start angle=180, end angle=128.7, radius=0.8];
  \node[font=\small, anchor=east] at (2.75,0.45) {angle};
  \eqlabl{(tb)}{(-0.6,2.6)}{torch, 30\,cm from the paper}
  \eqlab{(4.6,-0.05)}{(5.6,-0.8)}{light patch on 1\,cm grid paper}
\end{tikzpicture}
\expfigcap{Keep the torch the same distance from the paper at every angle.}
\end{center}

\subsubsection*{Equipment}
Torch (LED, with a narrow beam) \textbullet{} retort stand and clamp \textbullet{} 1\,cm grid paper \textbullet{} protractor \textbullet{} metre ruler \textbullet{} pencil \textbullet{} a darkened room

\needspace{10\baselineskip}
\task Complete the risk assessment. \skr{PC3}
\begin{ctab}
\begin{tabular}{|SW{4.0cm}|W{5.2cm}|W{5.4cm}|}
  \hline\hd & \thead{Risk (what could go wrong)} & \thead{Prevention} & \thead{Responding (if it happens)} \\ \hline
  & \scaf{Tripping in the darkened room} & \scaf{Bags under benches; walk, do not run} & \scaf{Tell the teacher; turn the lights on} \\ \hline
  & The retort stand tips over & & \\ \hline
  & Looking straight into the torch & & \\ \hline
\end{tabular}
\end{ctab}

\subsubsection*{Method}
\begin{enumerate}[itemsep=2pt]
  \item Tape the grid paper flat on the bench. Clamp the torch so the middle of its beam hits the centre of the paper at 90\textdegree, with the torch 30\,cm from the paper.
  \item Darken the room. Trace the edge of the light patch in pencil.
  \item Count the squares (the more-than-half rule) and record the area. A partner counts it again.
  \item Repeat with the beam at 75\textdegree, 60\textdegree, 45\textdegree, 30\textdegree{} and 15\textdegree{} to the paper, measuring the angle with the protractor and keeping the torch 30\,cm away.
  \item Do the whole set three times on fresh paper, then calculate the averages.
\end{enumerate}

% ─── RESULTS ─────────────────────────────────────────────────────────────────
\exppart{Results}{Amber}

\subsubsection*{Table 1.1: Area of the light patch at each angle \skr{PA1}}
\begin{ctab}
\begin{tabular}{|SW{3.0cm}|C{2.7cm}|C{2.7cm}|C{2.7cm}|C{3.0cm}|}
  \hline\hd & \thead{Angle of beam (\textdegree)} & \thead{Trial 1 (cm$^2$)} & \thead{Trial 2 (cm$^2$)} & \thead{Trial 3 (cm$^2$)} & \thead{Average (cm$^2$)} \\ \hline
  & \scaf{Example: 50} & \scaf{31} & \scaf{33} & \scaf{32} & \scaf{32} \\ \hline
  & 90 & & & & \\ \hline
  & 75 & & & & \\ \hline
  & 60 & & & & \\ \hline
  & 45 & & & & \\ \hline
  & 30 & & & & \\ \hline
  & 15 & & & & \\ \hline
\end{tabular}
\end{ctab}

\needspace{22\baselineskip}
\task Draw a line graph of the average area against the angle. Put the angle on the horizontal axis, give the graph a title, and draw a smooth line of best fit. \skr{PA2}
\graphgrid{10cm}

% ─── ANALYSIS ────────────────────────────────────────────────────────────────
\exppart{Analysis and evaluation}{Coral}

\task Describe the trend in your results. Use numbers from Table 1.1. \skr{PA5}
\answerbox{4cm}

\needspace{12\baselineskip}
\task The torch gives out the same amount of light at every angle. Explain what happens to the light on each square of paper when the patch gets bigger. Use the word \emph{intensity}.
\answerbox{4.5cm}

\needspace{14\baselineskip}
\task At midday in Melbourne the Sun is about 75\textdegree{} above the horizon in December and about 28\textdegree{} in June. Use your graph to explain why the ground heats up more in December. \skr{EV2}
\answerbox{5cm}

\needspace{10\baselineskip}
\task A student says, ``It's summer because Earth is closer to the Sun.'' Use your results and the idea of axial tilt to explain why this is wrong.
\answerbox{4.5cm}

\needspace{12\baselineskip}
\task Give two sources of error and how each could be reduced. \skr{EV3}
\begin{ctab}
\begin{tabular}{|SW{7.2cm}|W{7.8cm}|}
  \hline\hd & \thead{Source of error} & \thead{How to reduce it} \\ \hline
  & \scaf{The edge of the patch is blurry} & \scaf{Use a torch with a sharper beam; agree on the edge before tracing} \\ \hline
  & & \\ \hline
  & & \\ \hline
\end{tabular}
\end{ctab}

\needspace{10\baselineskip}
\task Write a conclusion: the trend, the numbers that back it, and whether your hypothesis was supported. \skr{EV1}
\answerbox{4.5cm}

\begin{yr8box}
\begin{keepext}
\exttask Relative brightness $=$ area at 90\textdegree{} $\div$ area at the angle. Calculate it for 30\textdegree{} from your averages. Then calculate $\sin 30^\circ$ on a calculator and compare. What does this suggest about the shape of your graph?
\workbox{5cm}
\end{keepext}
\end{yr8box}

% ─── SKILL CHECK ─────────────────────────────────────────────────────────────
\exppart{Skill check}{NavyBlue}

\begin{skillcheck}
\scrow{QP1}{I can write a hypothesis that names both variables.}
\scrow{PC2}{I changed only the angle and kept the distance the same.}
\scrow{PC4}{I counted each patch with the same rule, three times.}
\scrow{PA1}{My table has a title, headings, units and matching decimals.}
\scrow{PA2}{I drew a smooth line of best fit, not dot-to-dot.}
\scrow{PA5}{I can describe the trend, quoting numbers.}
\scrow{EV2}{I can use my results to explain the seasons.}
\end{skillcheck}

\task Which skill will you focus on in Experiment 2? Explain what you will do differently.
\answerbox{5cm}
